\section{每天怎么复盘：一张表就够}

收盘后把候选分成三栏：项目直接关系层、平台页面敞口层和广义AI内容扩散层。第一栏才有资格竞争主板龙头；第二栏用来判断热度；第三栏只看扩散，不能替第一栏背书。

\begin{exhibitbox}[展品：收盘检查表]
\begin{tabularx}{\textwidth}{L{2.3cm}X L{3.0cm}}
\toprule
收盘时记录什么 & 合格标准 & 得出的动作 \\
\midrule
新催化 & 有时间戳的一手资料，且提供播出反馈或第二项目。 & 继续看D+2。 \\
主板高度 & 项目直接关系主板股从首板走到二板，并有同层首板。 & 高度打法进入候选。 \\
主板容量 & 项目直接关系主板股连续两日成交额前二，收盘不弱于前收。 & 容量打法进入候选。 \\
日终承接 & 收盘不弱于前收，且未全天一字锁住。 & 次日保留观察资格。 \\
板块广度 & 合格集合至少三只，三分之二收涨。 & 题材继续保留。 \\
\bottomrule
\end{tabularx}
\end{exhibitbox}

纸面交易动作也可以写得很简单：D0只建观察表；D+1只确认二板和承接；D+2才决定题材是否继续；D+3到D+5只在第二催化到来时重看强化机会。盘中不追一字，不把开盘冲高当承接，不用传闻补公司关系。

掌阅科技、博纳影业和中广天择当前位于广义AI内容扩散层。它们走强时，结论只能是“AI内容在扩散”；主板项目龙头的结论仍要等项目直接关系公司出现。\sourcenote{analysis/non\_chinext\_candidate\_audit.md；S06；S09。}
